\documentclass[10pt,a4paper]{amsart}
\usepackage{amsmath,amssymb,amsthm,graphicx}
\usepackage{verbatim}
\newenvironment{reference}{\comment}{\endcomment}
%\newenvironment{reference}{}{}
\newenvironment{slogan}{\comment}{\endcomment}
\newenvironment{history}{\comment}{\endcomment}

% For commutative diagrams we use Xy-pic
\usepackage[all]{xy}

% We use 2cell for 2-commutative diagrams.
\xyoption{2cell}
% Theorem environments.
%
\theoremstyle{plain}
\newtheorem{theorem}[subsection]{Theorem}
\newtheorem{proposition}[subsection]{Proposition}
\newtheorem{lemma}[subsection]{Lemma}

\theoremstyle{definition}
\newtheorem{definition}[subsection]{Definition}
\newtheorem{example}[subsection]{Example}
\newtheorem{exercise}[subsection]{Exercise}
\newtheorem{situation}[subsection]{Situation}

\theoremstyle{remark}
\newtheorem{remark}[subsection]{Remark}
\newtheorem{remarks}[subsection]{Remarks}

\numberwithin{equation}{subsection}

% Macros
%
\def\lim{\mathop{\mathrm{lim}}\nolimits}
\def\colim{\mathop{\mathrm{colim}}\nolimits}
\def\Spec{\mathop{\mathrm{Spec}}}
\def\Hom{\mathop{\mathrm{Hom}}\nolimits}
\def\Ext{\mathop{\mathrm{Ext}}\nolimits}
\def\SheafHom{\mathop{\mathcal{H}\!\mathit{om}}\nolimits}
\def\SheafExt{\mathop{\mathcal{E}\!\mathit{xt}}\nolimits}
\def\Sch{\mathit{Sch}}
\def\Mor{\mathop{\mathrm{Mor}}\nolimits}
\def\Ob{\mathop{\mathrm{Ob}}\nolimits}
\def\Sh{\mathop{\mathit{Sh}}\nolimits}
\def\NL{\mathop{N\!L}\nolimits}
\def\CH{\mathop{\mathrm{CH}}\nolimits}
\def\proetale{{pro\text{-}\acute{e}tale}}
\def\etale{{\acute{e}tale}}
\def\QCoh{\mathit{QCoh}}
\def\Ker{\mathop{\mathrm{Ker}}}
\def\Im{\mathop{\mathrm{Im}}}
\def\Coker{\mathop{\mathrm{Coker}}}
\def\Coim{\mathop{\mathrm{Coim}}}

% Boxtimes
%
\DeclareMathSymbol{\boxtimes}{\mathbin}{AMSa}{"02}

%
% Macros for moduli stacks/spaces
%
\def\QCohstack{\mathcal{QC}\!\mathit{oh}}
\def\Cohstack{\mathcal{C}\!\mathit{oh}}
\def\Spacesstack{\mathcal{S}\!\mathit{paces}}
\def\Quotfunctor{\mathrm{Quot}}
\def\Hilbfunctor{\mathrm{Hilb}}
\def\Curvesstack{\mathcal{C}\!\mathit{urves}}
\def\Polarizedstack{\mathcal{P}\!\mathit{olarized}}
\def\Complexesstack{\mathcal{C}\!\mathit{omplexes}}
% \Pic is the operator that assigns to X its picard group, usage \Pic(X)
% \Picardstack_{X/B} denotes the Picard stack of X over B
% \Picardfunctor_{X/B} denotes the Picard functor of X over B
\def\Pic{\mathop{\mathrm{Pic}}\nolimits}
\def\Picardstack{\mathcal{P}\!\mathit{ic}}
\def\Picardfunctor{\mathrm{Pic}}
\def\Deformationcategory{\mathcal{D}\!\mathit{ef}}
\usepackage{cleveref}
\crefname{originallemma}{original lemma}{original lemmas}
\Crefname{originallemma}{Original lemma}{Original lemmas}
\makeatletter
\newlabel{formal-spaces-lemma-etale}{{0AN8}{source}{}{sourcepointer-1}{}}
\newlabel{formal-spaces-lemma-etale@cref}{{[originallemma][0][]0AN8}{[1][0][]source}{}{}{}}
\makeatother
\crefname{originallemma}{original lemma}{original lemmas}
\Crefname{originallemma}{Original lemma}{Original lemmas}
\makeatletter
\newlabel{formal-spaces-lemma-morphism-between-formal-spectra}{{0AN0}{source}{}{sourcepointer-2}{}}
\newlabel{formal-spaces-lemma-morphism-between-formal-spectra@cref}{{[originallemma][0][]0AN0}{[1][0][]source}{}{}{}}
\makeatother
\crefname{originallemma}{original lemma}{original lemmas}
\Crefname{originallemma}{Original lemma}{Original lemmas}
\makeatletter
\newlabel{more-algebra-lemma-Dedekind-etale-extension}{{0AP2}{source}{}{sourcepointer-3}{}}
\newlabel{more-algebra-lemma-Dedekind-etale-extension@cref}{{[originallemma][0][]0AP2}{[1][0][]source}{}{}{}}
\makeatother
\crefname{originallemma}{original lemma}{original lemmas}
\Crefname{originallemma}{Original lemma}{Original lemmas}
\makeatletter
\newlabel{more-algebra-lemma-completion-dvr}{{0AP1}{source}{}{sourcepointer-4}{}}
\newlabel{more-algebra-lemma-completion-dvr@cref}{{[originallemma][0][]0AP1}{[1][0][]source}{}{}{}}
\makeatother
\title{Etale surjective morphisms of locally Noetherian formal spaces are rig-surjective}
\author{The Stacks Project Authors}
\date{stacks-project:restricted.tex:local-snapshot-20261012; GNU FDL 1.2 or later}
\begin{document}
\maketitle
\noindent\textit{Editorial note.} All delivered mathematical text is the author's own, copied from the pinned Stacks Project chapter file restricted.tex; no statement or proof step was rewritten, reordered or removed.\par\medskip
\noindent\textit{Editorial note.} The counted proof is the author's complete argument: the proof base-changes the test object, extracts an affine etale piece, finds a maximal ideal over the closed point, promotes its local ring to a complete DVR, and composes the resulting formal spectrum morphisms\par\medskip
\noindent\textit{Editorial note.} The label lemma-etale-covering-rig-surjective names the counted statement, and the proof environment immediately following it in the source is delivered verbatim as its proof.\par\medskip
\noindent\textit{Editorial note.} The counted statement is carried by the author's own lemma environment in the chapter source.\par\medskip
\noindent\textit{Allowed standard prerequisites (editorial).}\begin{itemize}
\item rig-surjectivity tested by adic morphisms from complete discrete valuation rings
\item base change of representable etale morphisms and nonemptiness of a surjective etale cover
\item etale morphisms of formal algebraic spaces correspond to etale ring maps on formal spectra (Formal Spaces, Lemma formal-spaces-lemma-etale)
\item etale local rings over a DVR are discrete valuation rings (More on Algebra, Lemma more-algebra-lemma-Dedekind-etale-extension) and their completions are complete DVRs (More on Algebra, Lemma more-algebra-lemma-completion-dvr)
\item morphisms between formal spectra (Formal Spaces, Lemma formal-spaces-lemma-morphism-between-formal-spectra)
\end{itemize}\medskip
\ifdefined\hypertarget\hypertarget{sourcepointer-1}{}\fi
\noindent\textit{Source reference.} Original lemma 0AN8: sources/stacks-chow/tags.txt, source lines 13431--13431. Prerequisite pointer invoked at proof line(s) 5377; it names a standard fact of a sibling chapter that the counted argument uses, and every step of the counted proof built on it is retained verbatim in the delivered body.\par\medskip
\ifdefined\hypertarget\hypertarget{sourcepointer-2}{}\fi
\noindent\textit{Source reference.} Original lemma 0AN0: sources/stacks-chow/tags.txt, source lines 13423--13423. Prerequisite pointer invoked at proof line(s) 5385; it names a standard fact of a sibling chapter that the counted argument uses, and every step of the counted proof built on it is retained verbatim in the delivered body.\par\medskip
\ifdefined\hypertarget\hypertarget{sourcepointer-3}{}\fi
\noindent\textit{Source reference.} Original lemma 0AP2: sources/stacks-chow/tags.txt, source lines 13460--13460. Prerequisite pointer invoked at proof line(s) 5381; it names a standard fact of a sibling chapter that the counted argument uses, and every step of the counted proof built on it is retained verbatim in the delivered body.\par\medskip
\ifdefined\hypertarget\hypertarget{sourcepointer-4}{}\fi
\noindent\textit{Source reference.} Original lemma 0AP1: sources/stacks-chow/tags.txt, source lines 13459--13459. Prerequisite pointer invoked at proof line(s) 5383; it names a standard fact of a sibling chapter that the counted argument uses, and every step of the counted proof built on it is retained verbatim in the delivered body.\par\medskip
\begin{lemma}
\label{lemma-etale-covering-rig-surjective}
Let $S$ be a scheme. Let $f : X \to Y$ be a morphism of locally Noetherian
formal algebraic spaces which is representable by algebraic spaces, \'etale,
and surjective. Then $f$ is rig-surjective.
\end{lemma}

\begin{proof}
Let $p : \text{Spf}(R) \to Y$ be an adic morphism where $R$ is a complete
discrete valuation ring. Let $Z = \text{Spf}(R) \times_Y X$. Then
$Z \to \text{Spf}(R)$ is representable by algebraic spaces, \'etale, and
surjective. Hence $Z$ is nonempty. Pick a nonempty affine formal algebraic
space $V$ and an \'etale morphism $V \to Z$ (possible by our definitions).
Then $V \to \text{Spf}(R)$ corresponds to $R \to A^\wedge$ where
$R \to A$ is an \'etale ring map, see Formal Spaces, Lemma
\ref{formal-spaces-lemma-etale}. Since $A^\wedge \not = 0$
(as $V \not = \emptyset$) we can find a maximal ideal $\mathfrak m$
of $A$ lying over $\mathfrak m_R$. Then $A_\mathfrak m$ is a discrete
valuation ring (More on Algebra, Lemma
\ref{more-algebra-lemma-Dedekind-etale-extension}).
Then $R' = A_\mathfrak m^\wedge$ is a complete discrete valuation ring
(More on Algebra, Lemma \ref{more-algebra-lemma-completion-dvr}).
Applying Formal Spaces, Lemma
\ref{formal-spaces-lemma-morphism-between-formal-spectra}.
we find the desired morphism $\text{Spf}(R') \to V \to Z \to X$.
\end{proof}

\end{document}
